\documentclass[11pt]{article}

\usepackage[margin=1in]{geometry}
\usepackage[T1]{fontenc}
\usepackage[utf8]{inputenc}
\usepackage{lmodern}
\usepackage{microtype}
\usepackage{amsmath,amssymb,amsthm}
\usepackage{enumitem}
\usepackage{booktabs}
\usepackage{xurl}
\usepackage[hidelinks]{hyperref}

\newtheorem{definition}{Definition}
\newtheorem{lemma}{Lemma}
\newtheorem{remark}{Remark}

\title{Routing-Signature Compression for Many-Testing-Safe Anonymous DHT Tuning\\
\large Replacing $M$ by an effective multiplicity $M_{\mathrm{eff}}$}
\author{}
\date{March 21, 2026 (archive v493)}

\begin{document}
\maketitle

\begin{abstract}
Anonymous DHT tuning exposes a combinatorial explosion of knob settings.
Naively treating each knob vector as a distinct hypothesis can make certified menus overly conservative.

We show that in many anonymous routing designs, a large policy pool collapses to a much smaller set of \emph{distinct observable behaviors} on the support of contexts that actually occur.
We formalize this via \emph{routing signatures} and give a safe compression protocol that replaces the raw pool size $M$ by an effective multiplicity $M_{\mathrm{eff}}$ inside family-wise error control, without weakening the guarantee.
The result is an intentionally narrow optional companion to Certified~A rather than a second certified-routing entrypoint: Certified~A already yields a complete public menu claim with raw $M$, while the present note only explains when that accounting can be compressed safely.
This paper therefore stops at a compact \emph{signature manifest}: declared transcript projection, context partition, merge tolerance, quotient map, and resulting $M_{\mathrm{eff}}$.
Certified~A consumes that object when forming menus; Certified~C handles drift/change-control; Eval~1 and ABOM/OINL~+~Release~A handle notarization and public packaging.
A small DHT-style cutover drill and a remanifest/recertify ladder make that stop public-facing: later terse papers can now quote one exact verdict about when a fixed manifest still licenses compression, when it must be re-emitted under the same theorem root, and when the question has widened into Certified~C rather than speaking loosely about ``still basically the same menu."
\end{abstract}

\paragraph{Scope.}
We isolate and formalize \emph{Routing-Signature Compression for Many-Testing-Safe Anonymous DHT Tuning} as a reusable contract surface in the anonymity / Anonymous-DHT series. The paper is intentionally positioned as an optional certified-routing support contract beneath Certified~A rather than as an independent theorem root.

\paragraph{Imports.}
This paper owns only the multiplicity-compression step for Certified~A.
For the baseline many-testing story and menu guarantee, see Certified~A; for recertification under drift, see Certified~C.\cite{series:certifiedmenus,series:recert}
Signature-compression outputs are contract objects (Synthesis~0) and should be surfaced through ABOM/OINL manifests (Anonymity~C), while evaluator evidence release and signed public packaging belong to Eval~1 and Release~A.\cite{series:contractobjects,series:abomoinl,series:eval1,series:releaseA}

\paragraph{Exports.}
A routing-signature manifest: declared transcript projection, context partition, merge tolerance, representative-policy table, quotient map from raw knobs to signature classes, and the resulting effective multiplicity $M_{\mathrm{eff}}$. The public-facing claim is exactly that this manifest lets Certified~A replace raw-$M$ multiplicity accounting by a sound compressed count, not that it supersedes Certified~A's menu theorem.

\paragraph{Boundary.}
This paper stops once a fixed witness/interface has been quotiented into signature classes and $M_{\mathrm{eff}}$ has been computed.
Simultaneous menu selection belongs to Certified~A, disagreement-gated drift/change-control belongs to Certified~C, evaluator notarization belongs to Eval~1, and public lineage / release packaging belong to ABOM/OINL plus Release~A.\cite{series:certifiedmenus,series:recert,series:eval1,series:releaseA}

\paragraph{Earliest sufficient stop.}
If a verifier can inspect the declared transcript projection, context partition, merge tolerance, quotient map, representative-policy table, and the resulting $M_{\mathrm{eff}}$, then the only extra claim owned by this paper is already fixed. The menu theorem itself remains in Certified~A, and later drift, notarization, and release-lineage layers remain deliberately out of scope.

\paragraph{Later-terse-paper citation posture.}
Later terse papers should cite this note only when the live issue is the multiplicity-compression witness itself: which transcript projection, which context partition, which merge tolerance, and why raw $M$ collapsed to $M_{\mathrm{eff}}$. If the live issue is the baseline menu theorem, stop at Certified~A; if the live issue is drift / change-control for a deployed manifest, stop at Certified~C; and if the live issue is evaluator release or signed public lineage, stop at Eval~1 or Release~A rather than widening this note into a second certified-routing dossier.\cite{series:certifiedmenus,series:recert,series:eval1,series:releaseA}

\paragraph{Artifact policy.}
This archive is source-only; supporting artifacts are available on request. See \cite{series:artifactpolicy}.

\section{Signatures}
Policies differ in many internal details, but the adversary (and the privacy label) only depend on a chosen transcript projection $X$.

\begin{definition}[Routing signature]
Fix a context space $\mathcal C$ (e.g., local load bucket, churn bucket, epoch index) and a transcript projection $X$.
The \emph{signature} of a policy $\pi$ is the conditional law family $\{\mathcal L_\pi(X\mid C=c): c\in\mathcal C\}$.
Two policies are signature-equivalent if these conditional laws match on the support of observed contexts.
\end{definition}

\begin{remark}[Why signatures reduce multiplicity]
If two policies are signature-equivalent, then any statistical test or confidence bound that depends only on $X$ (and $C$) produces identical outcomes for the two policies.
Counting both in a multiplicity correction is pure waste.
\end{remark}

\section{Compression protocol}
Given randomized logs (or a simulator), compute a finite approximation of signatures by hashing conditional histograms.
A safe protocol is:
\begin{enumerate}[leftmargin=*]
\item Choose a finite context partition $\mathcal C=\bigsqcup_j \mathcal C_j$ and discretize $X$.
\item For each policy $\pi_m$ and each bin $\mathcal C_j$, estimate the histogram of $X$.
\item Define a signature string by concatenating these histograms (with a fixed ordering and rounding rule) and hashing.
\item Let $M_{\mathrm{eff}}$ be the number of distinct signature hashes.
\item Emit a \emph{signature manifest}: witness/interface version, context partition, rounding rule, merge tolerance, signature-class representatives, and the raw-to-compressed quotient map.
\item Run Certified~A's menu construction with $M$ replaced by $M_{\mathrm{eff}}$ in the $\alpha$-splitting.
\end{enumerate}

\section{Why this is safe}
\begin{lemma}[Multiplicity depends on distinct tests]
Any family-wise error control argument that uses a union bound (or a closed testing procedure) can replace a multiset of identical tests by a single representative without increasing error.
\end{lemma}

The only subtlety is approximation error: histogram discretization can merge distinct signatures.
This is handled by including a \emph{verification margin}: only merge signatures if they are within a certified tolerance in TV (or if the discretization is exact by construction).

\section{Exported object and earliest sufficient stop}
The public-facing output of this note is not a release receipt and not a certified menu.
It is the narrower \emph{signature manifest} that lets downstream papers say exactly why $M$ was replaced by $M_{\mathrm{eff}}$:
\begin{enumerate}[leftmargin=*]
\item the declared transcript projection and witness/interface version,
\item the context partition and discretization rule,
\item the merge tolerance (or exact-equivalence claim),
\item the quotient map from raw policy IDs to signature classes,
\item one representative policy ID per class, and
\item the resulting effective multiplicity $M_{\mathrm{eff}}$.
\end{enumerate}
If a verifier only needs to understand the many-testing correction in Certified~A, stopping at this manifest is sufficient.
The attempt log, anti-grinding evidence, and signed public lineage are intentionally deferred to Eval~1 and Release~A.\cite{series:eval1,series:releaseA}

\section{A tiny DHT-style example}
Suppose the declared transcript projection $X$ is ``committee-contact count bucket'' and the context partition is just two bins: weekday and weekend epochs.
A raw certified-menu pool with $M=12$ knob settings can still collapse sharply if those settings induce only three distinct conditional transcript laws on that declared support.
For example, the quotient map may be
\[
\{\pi_1,\ldots,\pi_5\}\mapsto s_A,\qquad
\{\pi_6,\ldots,\pi_9\}\mapsto s_B,\qquad
\{\pi_{10},\pi_{11},\pi_{12}\}\mapsto s_C,
\]
with exact signature equality inside each class under the declared discretization.
Then $M_{\mathrm{eff}}=3$ even though the raw tuning pool still has size $12$.
If Certified~A uses a simple Bonferroni split at family-wise level $\alpha=0.05$, the per-signature allocation rises from $0.05/12\approx 0.00417$ to $0.05/3\approx 0.0167$ without weakening the guarantee, because the declared projection has only three distinct tests to protect.

\paragraph{A minimal public signature-manifest card.}
\begin{center}\small
\begin{tabular}{@{}ll@{}}
Transcript projection & committee-contact count bucket \\
Contexts & weekday / weekend epochs \\
Merge tolerance & exact under the declared discretization \\
Quotient map & $\{\pi_1,\ldots,\pi_5\}\mapsto s_A$, $\{\pi_6,\ldots,\pi_9\}\mapsto s_B$, $\{\pi_{10},\pi_{11},\pi_{12}\}\mapsto s_C$ \\
Representatives & $\pi_1$, $\pi_6$, $\pi_{10}$ \\
Multiplicity summary & raw $M=12$, compressed $M_{\mathrm{eff}}=3$ \\
Certified~A handoff & replace raw multiplicity by $3$ for this fixed witness/interface version \\
\end{tabular}
\end{center}
This is the smallest public packet later terse papers need when the live issue is only why multiplicity compression was sound for the declared projection. Drift monitoring, re-merging, and lineage are all owned elsewhere.

\paragraph{Tiny boundary-slice drill.}
Keep the same projection and quotient map, but suppose a later deployment week includes a declaredly relevant ``holiday'' epoch slice that was outside the old weekday/weekend support and that one class member in $s_B$ now shows a distinct committee-contact bucket law on that slice.
The old manifest does \emph{not} license compression there, because the public claim was only about equality on the declared support of observed contexts.
A later terse paper may still quote the old manifest for weekday/weekend epochs and fail closed outside that support, but it may not silently reuse $M_{\mathrm{eff}}=3$ for the widened support.
The exact cutover is therefore visible: either keep the old support claim unchanged, or emit a new signature manifest that explicitly includes the holiday slice.

\paragraph{Minimal remanifest / recertify ladder.}
The smallest quotable cutover after a signature manifest exists is:
\begin{center}\small
\begin{tabular}{@{}p{0.30\linewidth}p{0.23\linewidth}p{0.37\linewidth}@{}}
\toprule
Observed change & Smallest public verdict & Required widening \\
\midrule
Raw knob IDs are renamed or permuted, but the declared projection, context support, merge tolerance, and quotient classes are unchanged & Reuse the old signature manifest & None; Certified~A may still use the same $M_{\mathrm{eff}}$ \\
A class representative changes or the quotient map is re-emitted with the same classes and support & Emit a refreshed manifest under the same compression claim & No new theorem root; keep the same Certified~A handoff \\
The declared support widens, a previously out-of-support slice becomes relevant, or the merge tolerance / discretization rule changes & Old manifest no longer licenses compression on the widened claim & Re-manifest first; then let Certified~C decide whether the old certified menu survives, needs incremental evidence, or must fail closed into a new claim object \\
\bottomrule
\end{tabular}
\end{center}
This ladder keeps the present paper narrow.
It states the smallest public verdict about the compression object itself, and then hands drift/change-control to Certified~C rather than re-owning recertification here.\cite{series:recert}

\section{Takeaways}
\begin{itemize}[leftmargin=*]
\item \textbf{Multiplicity should reflect distinct observable behaviors:} if two knob settings induce the same transcript law on realized contexts, they should not both count against the family-wise error budget.
\item \textbf{Routing signatures are the quotient:} define signatures as conditional laws of the transcript projection $X$ given context $C$, then compress by hashing discretized conditional histograms.
\item \textbf{Safety requires a tolerance story:} only merge signatures when equality is exact by construction or when a certified TV tolerance is enforced; otherwise discretization can hide real differences.
\item \textbf{Stop at the signature manifest:} the stable exported object is the quotient map and its tolerance/accounting metadata, not a second certified-menu theorem, full release receipt, or notarized evaluation bundle.
\item \textbf{Make the support cut public:} a manifest licenses compression only on its declared projection, context support, and merge rule; once those widen, the old $M_{\mathrm{eff}}$ may still be cited on the old support but may not be silently reused on the widened claim.
\item \textbf{Where it plugs in:} Certified~A's menu construction replaces $M$ by $M_{\mathrm{eff}}$ after signature compression, while Certified~C / Eval~1 / Release~A own later drift, notarization, and packaging stages.
\end{itemize}

\begin{thebibliography}{99}\small
\bibitem{series:artifactpolicy}
Anonymous.
\newblock Artifact and Disclosure Policy for the One-True-Archive (Source-Only Public Release).
\newblock Series synthesis note (Synthesis~6), 2026. (This archive.)
\bibitem{series:certifiedmenus}
Anonymous.
\newblock Certified Menus for Anonymous DHT Lookups: Distribution-free parameter selection with simultaneous guarantees under many-testing.
\newblock Certified series Paper~A, 2026. (This archive.)

\bibitem{series:contractobjects}
Anonymous.
\newblock Contract Objects for Anonymous-DHT Privacy Budgets and Claims.
\newblock Series synthesis note (Synthesis~0), 2026. (This archive.)

\bibitem{series:abomoinl}
Anonymous.
\newblock ABOM/OINL: Audit BOM and Outcome Identification Noise Label for Anonymous DHT Deployments.
\newblock Anonymity series Paper~C, 2026. (This archive.)

\bibitem{series:recert}
Anonymous.
\newblock Disagreement-Gated Recertification and Boundary Audits for Anonymous DHT Routers.
\newblock Certified series Paper~C, 2026. (This archive.)

\bibitem{series:eval1}
Anonymous.
\newblock Notarized Anonymity Certificates: Attempt Logs, Spending Discipline, and Verifier Rules.
\newblock Evaluation series Paper~1, 2026. (This archive.)

\bibitem{series:releaseA}
Anonymous.
\newblock Release Property Ledgers and Signed Receipts for Anonymous-DHT Privacy Claims.
\newblock Release and destination series Paper~A, 2026. (This archive.)

\end{thebibliography}

\end{document}
